\question[4]
Beim Einrichten einer neuen App auf dem Smartphone fragt diese nach 
verschiedenen Berechtigungen, unter anderem nach Zugriff auf Kontakte, 
Kamera und Standort.

\begin{enumerate}
    \item Erkläre, was mit der ,,Rechteverwaltung`` (Berechtigungen) 
    von Apps gemeint ist und wie sie dabei hilft, das Sammeln 
    personenbezogener Daten einzuschränken.
    \item Eine Taschenlampen-App fordert bei der Installation Zugriff 
    auf die Kontakte und den Standort des Geräts an. Bewerte diese 
    Forderung und begründe, wie ein Nutzer in diesem Fall sinnvoll 
    vorgehen sollte.
    \item Nenne eine weitere Einstellungsmöglichkeit (außer der 
    App-Rechteverwaltung), mit der sich die Weitergabe des eigenen 
    Standorts einschränken lässt.
\end{enumerate}

\begin{solutionbox}{8cm}

\textbf{1.} Bei der Rechteverwaltung legt das Betriebssystem fest, 
dass Apps für den Zugriff auf bestimmte Funktionen oder Daten des 
Geräts (z.\,B.\ Kamera, Mikrofon, Kontakte, Standort) explizit die 
Erlaubnis der Nutzerin oder des Nutzers einholen müssen. Dadurch kann 
gezielt entschieden und jederzeit wieder geändert werden, welche App 
auf welche persönlichen Daten zugreifen darf, statt dass jede App 
automatisch uneingeschränkten Zugriff auf alle Daten des Geräts hat.

\textbf{2.} Der Zugriff auf Kontakte und Standort ist für die 
Funktion einer Taschenlampen-App nicht notwendig. Eine solche 
Berechtigungsanfrage ist daher als unangemessen zu bewerten und ein 
Hinweis darauf, dass die App möglicherweise personenbezogene Daten zu 
anderen Zwecken sammeln möchte (z.\,B.\ für Werbung oder Weiterverkauf 
der Daten). Ein Nutzer sollte diese Berechtigungen verweigern bzw.\ 
nicht erteilen; funktioniert die App danach nicht mehr wie erwartet 
oder verlangt sie die Rechte zwingend, sollte eine alternative App 
ohne diese Anforderungen gewählt werden.

\textbf{3.} Weitere Möglichkeiten, z.\,B.: Deaktivieren der 
Standortdienste (GPS) auf Betriebssystemebene, wenn sie nicht benötigt 
werden; Freigabe des Standorts nur ,,einmalig`` oder ,,nur während der 
Nutzung der App`` statt dauerhaft; Deaktivieren der WLAN- und 
Bluetooth-Standorterkennung.

\end{solutionbox}